% ============================================================
%  Section 8 — Skill optimization and measured prompt iteration
% ============================================================
\section{Company-Specific Skill Adaptation}
\label{sec:skillopt}

\subsection{SkillOpt Integration and Evaluation}

DocDesk represents company behaviour as versioned Markdown skills. The actual
SkillOpt 0.2.0 library runs against the Extractor environment through the
product API. It receives labeled TRAIN examples, groups failure feedback into
bounded skill patches, then evaluates candidate changes on a disjoint SELECTION
split. A candidate is promoted only when the configured continuous score
improves. Model weights are unchanged; TEST remains reserved for a final
comparison.

The Company B run used 12 TRAIN, 20 SELECTION, and 10 TEST documents, covering
invoices, purchase orders, and receipts. The dataset snapshots and SHA-256
hashes are saved with the run. Kimi-K2.5 served both the target and optimizer
roles. The initial skill scored 1.00 hard exact match on SELECTION. Across
three four-example TRAIN steps, SkillOpt skipped all edits because selection
performance did not improve; the initial skill remained the best version. The
run took 1,041.5 seconds and recorded 97,805 tokens across 37 calls. These
values are specific to this small generated pilot, not a general-purpose
accuracy estimate.

The first baseline TEST pass returned exact scores for 4 of 10 documents; six
were lost to temporary provider connection errors. Retrying those six produced
valid exact predictions for all six. A fresh run of the selected skill also
completed all 10 documents with exact predictions. Thus both the recovered
baseline and selected skill scored 10/10 exact on this held-out set, while the
selected skill is byte-identical to the initial skill. This is evidence of no
measured change, not an optimization gain. Retry-aware counts and per-document
records are in \url{eval/results/company_b_extractor_skillopt.json}; raw
outputs and dataset fingerprints are preserved under
\url{skillopt_output/company_b_extractor/runs/20261019T155925Z}.

The previously reported broad Extractor NEM is not comparable: its scorer
counted \texttt{reference\_unique}, which the production response schema does
not return, and failed to map gold \texttt{line\_descriptions} to nested
production \texttt{lines[].description}. Both mismatches guaranteed field
errors. The scorer is corrected in the repository, but the broad baseline has
not been rerun; its old Extractor NEM values are withdrawn. The SkillOpt pilot's
1.00 selection score uses the corrected field mapping on only 20 SELECTION items,
so it remains an adapter-level result rather than a broad production accuracy
claim. The optimizer made no skill change and provides no evidence of an
accuracy gain.

To make the adaptation mechanism visible, a separate demo-only run started from
an intentionally poor prompt that always returned \texttt{UNKNOWN}. On 20
SELECTION documents, its soft normalized exact-match score rose from 0.01 to
0.21 after SkillOpt accepted bounded edits. Hard exact match remained 0.00 for
both prompts. On the untouched 10-document TEST split, soft NEM rose from 0.02
to 0.17, but hard exact match remained 0/10 for both. This shows improved
partial-field overlap, not successful complete extraction. The weak prompt is
isolated from the production skill and is included only to demonstrate the
optimizer. Artifacts are stored under
\url{skillopt_output/company_b_extractor_showcase}.

\subsection{Separate Manual Prompt Iteration}

We used the existing Company B TRAIN labels to investigate a measured
Coder error. Seven training documents showed that TechVision consulting,
audit, and research work maps to \texttt{6226-HO001}, while technical IT work
maps to \texttt{6170-IT001}. The v003 prompt added this exception. The broader
v004 revision added an explicit supplier-to-account map, made supplier identity
the primary signal for this benchmark, preserved line-description
disambiguation for TechVision's mixed services, and required short verbatim
citations.

\begin{table}[ht]
\caption{Company B Coder prompt comparison on the SELECTION split}
\label{tab:skill_iteration}
\centering\scriptsize\renewcommand{\arraystretch}{1.2}
\begin{tabular}{@{}lrrr@{}}
\toprule
\textbf{Metric} & \textbf{Baseline v001} & \textbf{Improved v004} & \textbf{Change} \\
\midrule
Cases (paired SELECTION) & 10 & 10 & -- \\
JSON parse rate & 10/10 & 10/10 & 0 \\
Correct account code & 7/10 & 9/10 & +2 \\
Supported code and citation & 2/10 & 6/10 & +4 \\
\bottomrule
\end{tabular}
\end{table}

Both prompts received the same gold extracted fields and identical
company-filtered top-5 RAG results, isolating the coding skill from OCR and
extraction quality. On the paired sample, v004 improved code accuracy by 20
percentage points and the stricter supported score by 40 points; it won four
cases, lost none, and tied six. All 10 responses parsed. This remains a small
selection-only prompt experiment, not a full-corpus or held-out TEST estimate.
The earlier four-case TechVision comparison remains a development check. Full
per-document outputs are in
\url{eval/results/company_b_coder_improvement.json}; prompt text and notes are
in \texttt{prompts/PROMPT\_LIBRARY.md}.

A separate 30-document v001/v002 run was throttled by provider rate limits. Only
12 documents had valid responses from both prompts; their supported Coder
scores were 9/12 and 8/12, respectively. This incomplete run is not evidence
that v002 improved performance. It motivated bounded requests, checkpointing,
and explicit separation of provider errors from model errors in the evaluator.

\subsection{Interpretation}

The v004 result is a manual, train-informed prompt revision; it is separate
from SkillOpt and must not be described as optimizer-generated. It is early
paired evidence from 10 SELECTION examples. The production skill loader selects
v004; TEST evaluation remains necessary before claiming generalization. The
older targeted comparison is available in
\texttt{eval/results/company\_b\_skill\_adaptation.json}.
